% ---
% titre: Priorité des puissances
% theme: NO
% chapitre: Nombres relatifs
% sous_chapitre: Puissances et racines
% niveau: [11e]
% type: EVAL
% tags: [c-puissance,c-racine,p-question-TS]
% difficulte: 1
% corrige: true
% ---


\question[10]
Écris sous forme de puissance.

\begin{parts}
	\vspace{20pt}
	\part
		$5 + 5  + 5 + 5 + 5= $\sol{\;\bl}{$5 \cdot 5 \textbf{\,(0.5pt)} = 5^2 \textbf{\,(0.5pt)}$}
	
	\vspace{20pt}
	\part
		$8 \cdot 8 \cdot 8 \cdot 8 \cdot 8 \cdot 8 \cdot 8 \cdot 8 = $\sol{\;\bl}{$8^8 \textbf{\,(1pt)}$}
	
	\vspace{20pt}
	\part
		$(-2)^2 + (-2)^2 + (-2)^2 + (-2)^2 =$\sol{\;\bl}{$16 \textbf{\,(0.5pt)} = 4^2 \textbf{\,(0.5pt)}$}
	
	\vspace{20pt}
	\part
		$9^4 \div 9^2 =$\sol{\;\bl}{$9^2 \textbf{\,(1pt)}$}
				
	\vspace{20pt}
	\part
		$(7^4)^3 = $\sol{\;\bl}{$7^{12} \textbf{\,(1pt)}$}
				
	\vspace{20pt}
	\part
		$(5^3)^5 \cdot 5^7 =$\sol{\;\bl}{$5^{15} \cdot 5^7 \textbf{\,(0.5pt)} = 5^{105} \textbf{\,(0.5pt)}$}
	
	\vspace{20pt}
	\part
		$12^{15} : 12^{14} =$\sol{\;\bl}{$12^1 \textbf{\,(1pt)}$}

	\vspace{20pt}
	\part
		$6^4 \cdot 6^{-5} =$\sol{\;\bl}{$6^{-1} \textbf{\,(0.5pt)}$}
	
	\vspace{20pt}
	\part
		$2^6 + 6^2 =$\sol{\;\bl}{$64  \textbf{\,(0.5pt)}+ 36 \textbf{\,(0.5pt)}=  100 \textbf{\,(0.5pt)} = 10^2 \textbf{\,(0.5pt)}$}	
\end{parts}


